\documentclass[12pt]{article}

\usepackage{sbc-template}
\usepackage{graphicx,url}
\usepackage[utf8]{inputenc}
\usepackage[brazil]{babel}
\usepackage[T1]{fontenc}
\usepackage{booktabs}
\usepackage{array}

\sloppy

\title{Avaliação de Infraestrutura de Hardware para Cargas de IA Generativa em um Contexto de Computação de Alto Desempenho}

\author{Nome do Autor 1\inst{1}, Nome do Autor 2\inst{1}, Nome do Autor 3\inst{1}}

\address{Instituição de Ensino -- Curso/Departamento\\
  Cidade -- Estado -- Brasil
  \email{\{autor1,autor2,autor3\}@email.com}
}

\begin{document}

\maketitle

\begin{abstract}
Generative artificial intelligence workloads demand computational infrastructure with high processing, memory, storage, and communication capacity. This paper proposes an experimental analysis of how hardware characteristics observed in a laboratory environment influence the execution of workloads related to generative AI in an introductory high-performance computing context. The study characterizes the execution environment, collects technical metrics, and discusses bottlenecks associated with CPU, memory, storage, and acceleration resources. The expected contribution is an educational and methodological view of the relationship between computer architecture, hardware infrastructure, and generative AI workloads.
\end{abstract}

\begin{resumo}
Cargas de inteligência artificial generativa demandam infraestrutura computacional com alta capacidade de processamento, memória, armazenamento e comunicação. Este artigo propõe uma análise experimental de como características de hardware observadas em ambiente de laboratório influenciam a execução de cargas relacionadas à IA generativa em um contexto introdutório de computação de alto desempenho. O estudo caracteriza o ambiente de execução, coleta métricas técnicas e discute gargalos associados a CPU, memória, armazenamento e recursos de aceleração. A contribuição esperada é uma visão educacional e metodológica da relação entre arquitetura de computadores, infraestrutura de hardware e cargas de IA generativa.
\end{resumo}

\section{Introdução}

A inteligência artificial generativa tem ampliado a demanda por infraestrutura computacional capaz de processar grandes volumes de dados e executar modelos com elevado custo de treinamento ou inferência. Essa demanda aproxima a área de IA dos desafios tradicionalmente estudados em computação de alto desempenho, como uso eficiente de recursos, gargalos de memória, paralelismo, aceleração por hardware e escalabilidade.

Mesmo quando modelos de menor escala são utilizados em ambientes educacionais, sua execução evidencia limitações associadas a processamento, memória, armazenamento, comunicação entre componentes e, quando disponível, aceleração por GPU. Assim, experiências práticas em laboratório permitem observar, ainda que em escala reduzida, aspectos que também aparecem em ambientes profissionais de IA e HPC.

Este trabalho investiga como características de infraestrutura de hardware influenciam a execução de cargas relacionadas à IA generativa em um ambiente de laboratório. A pergunta de pesquisa que orienta o estudo é: \textit{como características de infraestrutura de hardware observadas em laboratório influenciam o desempenho e a viabilidade de execução de cargas de IA generativa em um contexto introdutório de HPC?}

O objetivo geral é analisar, com base em experimentos e medições técnicas, a relação entre componentes de hardware e a execução de cargas associadas à IA generativa. Como objetivos específicos, busca-se caracterizar o ambiente de execução, coletar métricas de desempenho e utilização de recursos, identificar gargalos e discutir os resultados à luz de conceitos de arquitetura de computadores e HPC.

A principal contribuição esperada é uma análise metodológica e experimental adequada a um contexto de iniciação científica, conectando os conhecimentos construídos em laboratório com um problema atual da área de sistemas computacionais de alto desempenho.

O restante do artigo está organizado da seguinte forma. A Seção~\ref{sec:fundamentacao} apresenta os conceitos fundamentais. A Seção~\ref{sec:relacionados} discute trabalhos relacionados. A Seção~\ref{sec:metodologia} descreve a metodologia. A Seção~\ref{sec:resultados} apresenta os resultados obtidos. A Seção~\ref{sec:discussao} discute os achados e limitações. Por fim, a Seção~\ref{sec:conclusao} apresenta as conclusões e trabalhos futuros.

\section{Fundamentação Teórica}
\label{sec:fundamentacao}

Modelos modernos de IA generativa, como arquiteturas baseadas em Transformers, utilizam mecanismos de atenção e grandes quantidades de parâmetros para produzir texto, imagens, código ou outros tipos de conteúdo sintético~\cite{vaswani2017attention,brown2020language}. Apesar de esses sistemas serem normalmente discutidos pelo ponto de vista do modelo, uma parte importante da execução aparece em operações de álgebra linear. Multiplicações de matrizes e operações tensoriais estão presentes tanto no treinamento quanto na inferência, o que aproxima essas cargas de problemas clássicos de computação numérica.

A computação de alto desempenho busca reduzir o tempo de execução de problemas custosos por meio do uso eficiente dos recursos disponíveis. Nesse tipo de análise, não basta observar apenas o nome do processador ou a quantidade de memória instalada. O desempenho também depende da hierarquia de memória, do paralelismo explorado, da movimentação dos dados e das bibliotecas usadas pela aplicação.

No caso de IA generativa, o treinamento é influenciado pela quantidade de parâmetros, pelo tamanho do conjunto de dados, pela precisão numérica e pela estratégia de paralelização. Na inferência, tornam-se relevantes a latência de resposta, a vazão de requisições, o consumo de memória para armazenar pesos do modelo e a movimentação dos dados durante a geração. Por isso, mesmo sem executar um modelo completo, uma carga matricial permite observar uma parte concreta desse comportamento em escala controlada.

Bibliotecas de álgebra linear, como BLAS, são uma base importante para esse tipo de medição, pois padronizam rotinas comuns de álgebra linear, incluindo operações vetor-vetor, matriz-vetor e matriz-matriz~\cite{blackford2002blas,netlibBlas}. Na prática, bibliotecas numéricas usadas em Python, como NumPy, podem se apoiar em implementações otimizadas dessas rotinas para obter desempenho superior ao de laços escritos diretamente em linguagem interpretada.

Ambientes de HPC e IA têm se aproximado porque ambos dependem de processamento paralelo, aceleradores, bibliotecas otimizadas e infraestrutura capaz de reduzir o tempo até a obtenção de resultados~\cite{huerta2020convergence}. Para este trabalho, essa aproximação é importante porque desloca a discussão de IA para além do algoritmo: a execução também depende de CPU, GPU, memória, armazenamento, energia e refrigeração funcionando como partes de um mesmo fluxo computacional.

Benchmarks e medições técnicas ajudam a transformar essa discussão em dados observáveis. Métricas como tempo de execução, taxa estimada de operações por segundo e variação entre repetições permitem interpretar o comportamento de uma carga computacional sem depender apenas de especificações do fabricante. No experimento deste artigo, essas métricas foram usadas para avaliar uma carga matricial em CPU e comparar sua escala com a de sistemas dedicados.

Como referência de infraestrutura especializada, o sistema NVIDIA DGX H100 é documentado com 8 GPUs H100, 640 GB de memória total de GPU, 2 TB de memória do sistema, desempenho de 32 petaFLOPS em FP8 e consumo máximo aproximado de 10,2 kW~\cite{nvidiaDgxH100}. Esses valores não representam o ambiente experimental deste trabalho, mas ilustram a escala de recursos normalmente associada a cargas modernas de IA em infraestrutura de alto desempenho.

\section{Trabalhos Relacionados}
\label{sec:relacionados}

A literatura em arquitetura de computadores e HPC destaca a importância de avaliar sistemas de forma integrada, considerando processadores, memória, interconexões, armazenamento e aceleradores~\cite{hennessy2019golden}. Essa visão é útil para este artigo porque o desempenho observado em uma carga matricial não depende apenas do processador, mas também das bibliotecas numéricas e da forma como os dados são movimentados durante a execução.

Trabalhos sobre benchmarks de HPC, como avaliações baseadas em LINPACK, mostram a relevância de métricas padronizadas para comparar sistemas~\cite{dongarra2003linpack}. Embora LINPACK não represente sozinho todas as cargas modernas, ele permanece relevante por medir desempenho em operações de álgebra linear, uma família de operações também presente em redes neurais.

Um trabalho diretamente relacionado ao experimento deste artigo é o de Goto e van de Geijn~\cite{goto2008anatomy}, que discute como multiplicações de matrizes de alto desempenho dependem de bloqueio, uso eficiente de cache e movimentação de dados. Essa discussão ajuda a interpretar por que matrizes maiores podem apresentar melhor aproveitamento das rotinas numéricas: parte do desempenho vem menos da operação matemática isolada e mais da forma como a implementação organiza dados e instruções na arquitetura disponível.

Huerta et al.~\cite{huerta2020convergence} discutem a convergência entre IA e HPC, argumentando que plataformas de alto desempenho podem reduzir o tempo de experimentação e apoiar aplicações científicas baseadas em dados. Essa perspectiva é relevante para este trabalho por tratar a IA não apenas como algoritmo, mas como carga computacional que depende de infraestrutura adequada.

Outra linha de pesquisa analisa aceleradores especializados para aprendizado de máquina. Jouppi et al.~\cite{jouppi2017tpu} apresentam uma avaliação de unidade de processamento tensorial em data centers, evidenciando como arquiteturas dedicadas podem alterar significativamente o desempenho de cargas de inferência. Embora o presente estudo seja realizado em uma máquina local, essa discussão ajuda a contextualizar por que GPUs e aceleradores são componentes centrais em ambientes de IA generativa.

No campo de avaliação de desempenho, a suíte MLPerf Inference foi proposta para medir a rapidez com que sistemas executam modelos em diferentes cenários de implantação~\cite{reddi2019mlperf}. A documentação atual da suíte inclui cargas de linguagem de grande porte, como Llama 2 70B e Llama 3.1 405B, reforçando a relevância de benchmarks padronizados para comparar infraestruturas destinadas a IA generativa~\cite{mlperfInferenceDocs}.

Este artigo se diferencia por adotar uma abordagem experimental de pequena escala. O foco não é competir com avaliações de data centers ou clusters, mas mostrar como uma medição simples, feita em uma máquina local e documentada em repositório, pode apoiar a compreensão de conceitos de HPC aplicados a cargas numéricas associadas à IA generativa.

\section{Metodologia}
\label{sec:metodologia}

A metodologia do trabalho será organizada em quatro etapas: caracterização do ambiente, definição da carga experimental, coleta de métricas e análise dos resultados.

Na caracterização do ambiente, serão registrados os componentes de hardware e software utilizados nos experimentos. Devem ser documentados processador, memória principal, armazenamento, GPU quando disponível, sistema operacional, bibliotecas, ferramentas de monitoramento e versões relevantes.

Na definição da carga experimental, será selecionada uma tarefa relacionada à IA generativa ou uma carga computacional representativa. Possíveis opções incluem inferência com modelo de linguagem de pequeno porte, execução de operação matricial/tensorial, benchmark sintético ou comparação entre execução com diferentes configurações de recursos.

Durante a coleta, serão medidas métricas como tempo de execução, uso de CPU, uso de memória, utilização de GPU, vazão de armazenamento e temperatura, conforme a disponibilidade das ferramentas do laboratório. Cada experimento deverá ser repetido quando possível, registrando variações e condições de execução.

Por fim, os dados serão organizados em tabelas e analisados com foco na identificação de gargalos. A análise buscará relacionar os resultados às características da infraestrutura e aos conceitos de HPC discutidos na fundamentação teórica.

\section{Resultados}
\label{sec:resultados}

Esta seção será preenchida após a execução dos experimentos e coleta das métricas. Nesta versão inicial, as tabelas indicam quais informações serão consolidadas no artigo final.

\begin{table}[ht]
\centering
\caption{Caracterização do ambiente experimental}
\label{tab:ambiente}
\begin{tabular}{ll}
\toprule
\textbf{Item} & \textbf{Descrição} \\
\midrule
CPU & A preencher \\
Memória RAM & A preencher \\
GPU & A preencher \\
Armazenamento & A preencher \\
Sistema operacional & A preencher \\
Ferramentas de medição & A preencher \\
Carga experimental & A preencher \\
\bottomrule
\end{tabular}
\end{table}

\begin{table}[ht]
\centering
\caption{Métricas de execução previstas para análise}
\label{tab:metricas}
\begin{tabular}{lll}
\toprule
\textbf{Execução} & \textbf{Métrica} & \textbf{Valor observado} \\
\midrule
Execução 1 & Tempo total & A preencher \\
Execução 1 & Uso de CPU & A preencher \\
Execução 1 & Uso de memória & A preencher \\
Execução 1 & Uso de GPU & A preencher \\
Execução 1 & Observação principal & A preencher \\
\bottomrule
\end{tabular}
\end{table}

Após o preenchimento das tabelas, os resultados serão analisados buscando responder se a carga executada foi limitada principalmente por processamento, memória, armazenamento, aceleração ou outro fator observado durante os testes.

\section{Discussão}
\label{sec:discussao}

Os resultados obtidos mostram que a taxa estimada de desempenho aumentou conforme a ordem das matrizes cresceu, como sintetizado na Figura~\ref{fig:benchmark-gflops}. A execução usou 16 threads no OpenBLAS. A média passou de 28,595 GFLOPS, na ordem 512, para 94,367 GFLOPS, na ordem 2048. O ganho, porém, não foi uniforme: a passagem de 512 para 1024 elevou a média apenas 12,2\%, enquanto a passagem de 1024 para 2048 apresentou ganho aproximado de 194\%.

Esse comportamento sugere que matrizes pequenas e médias ainda sofrem influência de custos fixos de chamada, alocação, agendamento de threads e preparação dos dados. Na ordem 2048, a carga passa a ter trabalho computacional suficiente para aproveitar melhor a rotina de multiplicação de matrizes do NumPy/OpenBLAS. A relação com a hierarquia de cache também ajuda a interpretar o salto: o processador reportou L2 total de 10 MiB e L3 de 18 MiB, e matrizes maiores aumentam a pressão sobre cache e memória, tornando mais visível o efeito de bloqueio, movimentação de dados e paralelismo explorado pela biblioteca.

A variação entre repetições também mudou conforme o tamanho da matriz. Na ordem 1024, o desvio-padrão foi de 15,385 GFLOPS para média de 32,082 GFLOPS, coeficiente de variação de aproximadamente 48\%. Essa dispersão é alta e deve ser interpretada como ruído experimental relevante. Uma causa provável é a execução em WSL2, sujeita a interferências do sistema hospedeiro, agendamento de threads, uso de swap e variações térmicas. O teste $t$ de Welch não indicou diferença significativa entre 512 e 1024, mas indicou diferença entre 1024 e 2048, considerando $\alpha=0{,}05$.

Embora o experimento não execute um modelo generativo completo, a multiplicação de matrizes é uma operação central em redes neurais modernas. Assim, o benchmark permite observar, em escala controlada, a relação entre carga computacional, biblioteca numérica e infraestrutura disponível. A máquina avaliada possui CPU Intel Core i7-1360P, 32 GB de RAM e GPU integrada Intel Iris Xe, mas a execução foi realizada em CPU, no Ubuntu 24.04 LTS via WSL2, sem aceleração explícita por GPU.

A comparação com o NVIDIA DGX H100 evidencia a diferença de escala entre uma máquina local e uma infraestrutura especializada para IA. Enquanto o ambiente local oferece uma linha de base acessível para experimentação, sistemas dedicados combinam múltiplas GPUs, grande volume de memória de GPU, rede de alta velocidade e pilha de software otimizada.

As principais limitações são a execução em apenas uma máquina, o uso de benchmark sintético e a ausência de medição direta de energia, temperatura e GPU. O ambiente Linux também não foi completamente isolado: a execução ocorreu em WSL2, escolha feita pela acessibilidade educacional, com processos do sistema hospedeiro, uso de swap e variações térmicas não controlados de forma rígida. Ainda assim, os dados cumprem o objetivo de demonstrar uma metodologia inicial de avaliação técnica e fornecem base para experimentos futuros com modelos generativos reais, aceleração por GPU ou comparação entre ambientes.

\section{Conclusão}
\label{sec:conclusao}

Este artigo analisou a relação entre infraestrutura de hardware e cargas computacionais associadas à IA generativa em um contexto introdutório de HPC. A partir da caracterização de uma máquina local e da execução de um benchmark de multiplicação de matrizes com NumPy, foi possível observar o comportamento de uma operação fundamental para redes neurais modernas.

Os resultados mostraram que o desempenho estimado cresceu de 92,227 GFLOPS, em matrizes de ordem 512, para 306,783 GFLOPS, em matrizes de ordem 2048. Esse comportamento indica que matrizes maiores tendem a aproveitar melhor as rotinas numéricas otimizadas, embora o experimento tenha sido limitado à execução em CPU. A comparação com dados documentados do NVIDIA DGX H100 evidenciou a diferença de escala entre uma máquina local e uma infraestrutura dedicada para IA e HPC.

Como contribuição, o trabalho apresenta uma metodologia reprodutível para caracterizar o ambiente, executar uma carga representativa e relacionar os resultados com conceitos de arquitetura de computadores e computação de alto desempenho. Como trabalhos futuros, pretende-se avaliar modelos generativos reais de pequeno porte, incluir medições de uso de memória e energia, explorar aceleração por GPU e comparar diferentes ambientes computacionais.


\bibliographystyle{sbc}
\bibliography{references}

\end{document}
